%=====================================================================
\section{The Block Universe}
\label{ch:block}
%=====================================================================

This final chapter addresses the nature of time, initial conditions, and the ultimate structure of the universe in DRLT. The conclusion is radical but unavoidable: the universe is not a process. It is a matrix.

\subsection{The universe is $G$}

The Gram matrix $G_{ij} = \langle\psi_i|\psi_j\rangle$ for all $N$ points contains \emph{all} physical information:
\begin{itemize}
\item Geometry (distances, curvature) from $|G_{ij}|$,
\item Forces (gauge connections, field strengths) from $\arg(G_{ij})$,
\item Matter (deviations from vacuum) from $\det(G_h) > 0$,
\item Quantum mechanics ($\hbar_h$) from hinge areas,
\item Coupling constants from the Binet--Cauchy decomposition.
\end{itemize}

Nothing exists outside $G$. There is no background space in which $G$ is embedded, no external time in which $G$ evolves, no observer outside $G$ who reads it. $G$ is the universe, complete and self-contained.

\subsection{Constraint propagation: the block is forced}

The block universe is not a philosophical choice --- it is a mathematical consequence of $\text{rank}(G) \leq 5$.

Consider $N$ vertices with $\psi_i \in \mathbb{C}^5$. The Gram matrix $G$ has $N^2$ complex entries but only $5N$ real parameters (the components of the $\psi_i$). For $N \gg 5$, $G$ is massively over-determined: any single row of $G$ (the inner products of one vertex with all others) constrains all remaining $\psi_j$ up to a $\text{U}(5)$ gauge freedom.

Concretely: given a single reference state $\psi_0$ and the $N - 1$ inner products $\{G_{0j}\}_{j=1}^{N-1}$, the rank-5 constraint propagates to determine the \emph{entire} configuration $\{\psi_j\}$ (up to unitary equivalence). The number of free real parameters is:
\begin{equation}
\text{DOF} = 2d \cdot N - d^2 = 10N - 25
\end{equation}
where $d^2 = 25$ is subtracted for the $\text{U}(5)$ gauge freedom. For the physical content (the $d^2 = 25$ Wishart eigenvalues), only 25 real numbers encode all of physics.

Numerical verification confirms this: starting from a random $N$-vertex network, iterating the local evolution $U_i = e^{-iH_i/\hbar_{\text{eff},i}}$ converges to a $W$-invariant fixed point within $\sim 30$ iterations. The block universe is the unique self-consistent solution --- it does not evolve because there is nothing else it could be.

\subsection{Time as gradient}

There is no time variable in DRLT. The action $S = \sum_h A_h\,\delta_h$ contains no $t$. The Gram matrix $G$ is static.

What we experience as time is the \emph{gradient direction} of $\det(G_h)$ across the lattice:
\begin{equation}
\text{``time''} = \text{direction of decreasing } \det(G_h) = \text{direction of increasing alignment}.
\end{equation}

In the direction of increasing $\det(G_h)$: vertices are more random, $\hbar$ is larger, physics is more quantum --- this is the ``past'' (toward the Big Bang). In the direction of decreasing $\det(G_h)$: vertices align, $\hbar$ decreases, physics becomes classical --- this is the ``future'' (toward heat death).

The second law of thermodynamics is the statement that we read $G$ in the direction of increasing alignment (decreasing $\det(G_h)$, increasing entropy).

\subsection{Cosmic history as diagonalization}

The entire history of the universe is encoded in the structure of $G$:

\begin{center}
\begin{tabular}{lcll}
\hline
Region of $G$ & $\text{rank}$ & $\det(G_h)$ & Epoch \\
\hline
Random block & $N$ & large & Big Bang: $\text{SU}(N)$, all forces unified \\
Rank reduction & $N \to 5$ & decreasing & Swap annihilation: atoms pair and die \\
Rank-5 plateau & $5$ & moderate & Our era: $\text{SU}(3)\!\times\!\text{SU}(2)\!\times\!\text{U}(1)$ \\
Further alignment & $5 \to 1$ & small & Far future: $3 \to 2 \to 1 \to 0$ \\
Diagonal block & $1$ & $0$ & Heat death: $G_{ij} = 1\;\forall(i,j)$, vacuum \\
\hline
\end{tabular}
\end{center}

This is not a narrative; it is a \emph{description of different regions of the same matrix}. The Big Bang and heat death coexist in $G$, just as the first and last pages of a book coexist in the book.

\subsection{The atom death sequence}

The dimensional atoms annihilate in order of size (Chapter~\ref{ch:whyd5}). Larger atoms have stronger self-coupling and freeze faster:

\begin{equation}
\text{SU}(N) \;\xrightarrow{\text{fast}}\; \text{SU}(3) \times \text{SU}(2) \times \text{U}(1) \;\xrightarrow{\text{slow}}\; \text{U}(1) \;\xrightarrow{\text{very slow}}\; 1
\end{equation}

\begin{center}
\begin{tabular}{cccl}
\hline
Atom & $N_\text{eff}$ & Freezing rate & Status \\
\hline
3 (strong) & 1 & Fastest & Confined (frozen, locally active) \\
2 (weak) & 2 & Medium & Broken (EW transition complete) \\
1 (EM) & $\infty$ & Slowest & Active (last survivor) \\
0 & --- & --- & Vacuum (end state) \\
\hline
\end{tabular}
\end{center}

We live in the era where atom 3 is confined, atom 2 is broken, and atom 1 is still active. This is not a special time --- it is the \emph{only} era with interesting physics, because:
\begin{itemize}
\item Before the rank-5 plateau: too hot, no stable structures.
\item After atom 1 dies: too cold, no interactions.
\item On the plateau: chemistry, biology, observers.
\end{itemize}

\subsection{No initial conditions}

A block universe has no ``initial'' state. The matrix $G$ simply \emph{is}. The question ``what were the initial conditions?'' is like asking ``what is on page zero of a book?'' --- the question is malformed.

All physical quantities derived in this work --- coupling constants, masses, mixing angles --- are properties of the \emph{rank-5 plateau}, independent of what $G$ looks like in the random block or the diagonal block. They depend on $(n_S, n_T) = (3, 2)$ and $c = 2$, which are structural constants of the plateau, not initial conditions.

The only quantity that depends on ``where in the block'' is the cosmological constant $\Lambda \sim 1/N_H$, which encodes the size of the plateau relative to the total matrix. This is an \emph{address}, not a parameter.

\subsection{The cosmic address $\varepsilon_0$}

The ghost correction parameter $\varepsilon_0 \approx 0.0038$ (Chapter~\ref{ch:ghosts}) is not a free parameter. It encodes our position within the block:
\begin{equation}
\varepsilon_0 = f(N_H, d) \sim N_H^{-1/k}
\end{equation}
for some exponent $k$ determined by the geometry.

Different positions in $G$ correspond to different $\varepsilon_0$:
\begin{itemize}
\item Early universe ($\varepsilon_0$ large): large ghost corrections, less precise ``plateau values.''
\item Current epoch ($\varepsilon_0 \approx 0.0038$): small corrections, plateau values nearly exact.
\item Far future ($\varepsilon_0 \to 0$): corrections vanish, but forces also vanish --- nothing to measure.
\end{itemize}

The coupling constants \emph{approach} their geometric plateau values ($8, 30, 6\pi^2$) as $\varepsilon_0 \to 0$, but never reach them exactly because the forces themselves fade. \textbf{Perfect physics is unobservable}: by the time the universe is clean enough to display exact values, there is nothing left to observe them.

\subsection{The complete chain: zero to universe}

\begin{equation}
\boxed{
\text{Points} \;\xrightarrow{\text{Frobenius}}\; \mathbb{C}
\;\xrightarrow{\text{Atomic}}\; d = 5
\;\xrightarrow{\text{Causal}}\; 3\!+\!2\!+\!1
\;\xrightarrow{\Lambda^3}\; 25
\;\xrightarrow{\text{Basel}}\; \alpha_i
\;\xrightarrow{\text{Simplex}}\; \text{SM}
\;\xrightarrow{\text{Block}}\; \text{Universe}
}
\end{equation}

\begin{center}
\begin{tabular}{rlll}
\hline
Step & From & To & How \\
\hline
0 & Nothing & Points with relations & (minimal assumption) \\
1 & Relations & $\mathbb{C}$ & Frobenius + phase + det \\
2 & $\mathbb{C}^N$ & $d = 5$ & Atomic uniqueness + swap \\
3 & $\mathbb{C}^5$ & SU(3)$\times$SU(2)$\times$U(1), $c\!=\!2$ & Chirality + density \\
4 & Hinges & $1 + 12 + 12 = 25$ & Binet--Cauchy + $c$-weight \\
5 & Channels & $\alpha_3, \alpha_2, \alpha_1$ & Basel sum + Gram rank \\
6 & $\alpha_i + v_H$ & 26 SM parameters & Simplex geometry \\
7 & $\hbar_h = 0$ (vacuum) & Cosmology (9 observables) & Block universe \\
8 & $\text{Tr}(G) = N$ & Ghost Sum Rule ($\Sigma\delta = 0$) & Completeness \\
\hline
& \multicolumn{2}{l}{\textbf{Free parameters:}} & \textbf{0} \\
& \multicolumn{2}{l}{\textbf{Observational inputs:}} & \textbf{0} \\
& \multicolumn{2}{l}{\textbf{Predictions:}} & \textbf{35+} \\
& \multicolumn{2}{l}{\textbf{Error structure:}} & \textbf{predicted (sum rule)} \\
\hline
\end{tabular}
\end{center}

\subsection{The universal rank cascade}

The derivation chain above starts from ``points with relations'' and arrives at $\mathbb{C}^5$. But there is a deeper way to see this.

\textbf{Begin with nothing specific.} Take $N$ points with random pairwise complex values, forming an $N \times N$ Gram matrix $G$ of \emph{any} rank. $N$ can be enormous. No constraint is imposed --- not $\mathbb{C}^5$, not rank $\leq 5$, not even a specific dimension.

\textbf{Sort by $W$ and vary the resolution.} At high resolution (large $N_\text{eff}$), every individual $\psi$ matters: the effective rank is large, all degrees of freedom are active, all forces are unified. At medium resolution, outliers are averaged away: effective rank drops, symmetries begin to separate. At low resolution, almost all variation is washed out: effective rank approaches 1, the matrix looks nearly diagonal, and physics goes silent.

This is not a dynamical process. It is a \emph{mathematical property of large random matrices}: when viewed at different scales, any sufficiently large matrix naturally stratifies into rank bands.

\begin{center}
\begin{tabular}{ccl}
\hline
Effective rank & Gauge structure & Era \\
\hline
$N$ (full) & $\text{SU}(N)$ & Unification: all forces identical \\
$\vdots$ & $\vdots$ & Swap annihilation cascade \\
$5$ & $\text{SU}(3)\!\times\!\text{SU}(2)\!\times\!\text{U}(1)$ & Our era: structure, chemistry, life \\
$3$ & $\text{SU}(3)$ frozen, $\text{SU}(2)$ fading & Late universe \\
$1$ & $\text{U}(1) \to 1$ & Heat death: indistinguishable \\
\hline
\end{tabular}
\end{center}

The rank-5 band is not special. It is not selected, tuned, or designed. It exists in \emph{every} sufficiently large random matrix, with measure-theoretic certainty. This resolves three foundational puzzles simultaneously:

\begin{enumerate}
\item \textbf{The anthropic principle dissolves.} ``Why does the universe have the right constants for life?'' becomes ``Why does a random matrix have a rank-5 band?'' --- and the answer is: it always does. Observers can only exist in a band with enough structure for chemistry (rank $> 1$) but enough stability for persistence (rank not too high). Rank $\sim 5$ is where both conditions hold. This is not selection; it is measure theory.

\item \textbf{Fine-tuning dissolves.} The coupling constants, masses, and mixing angles are not ``tuned'' to their values. They are the \emph{mathematical properties} of the rank-5 band of any large random Hermitian matrix. $1/\alpha_1 = 6\pi^2$ is not a coincidence --- it is what the rank-5 band looks like, universally.

\item \textbf{Cosmic necessity dissolves.} ``Why is there something rather than nothing?'' becomes ``Does a random matrix have rank structure?'' --- and the answer is: yes, with probability 1. The rank cascade is not contingent. It is guaranteed.
\end{enumerate}

In this picture, $d = 5$ is not derived from atomic dimensions or swap elimination. Those are \emph{descriptions of why the rank-5 band has the structure it does} --- correct and useful, but not the reason we are here. We are here because we are in the rank-5 band, and the rank-5 band is here because every large matrix has one.

The parameter $\varepsilon_0 \approx 0.0038$ is our coordinate within the band: $\text{SU}(3)$ is mostly frozen, $\text{SU}(2)$ is freezing, $\text{U}(1)$ is still active. A different $\varepsilon_0$ corresponds to a different address --- not a different universe.

\subsection{Measurement is reading $G$}

In standard quantum mechanics, measurement requires a separate axiom (the Born rule or the projection postulate). In DRLT, no measurement axiom exists or is needed.

\textbf{Microscopic measurement} $=$ the inner product $G_{ij} = \langle\psi_i|\psi_j\rangle$ between two vertices. This is the same operation as spatial displacement (comparing two points) and temporal evolution (comparing two states). There is no distinction between ``interaction'' and ``measurement'' --- both are $G_{ij}$.

\textbf{Macroscopic ``collapse''} is a consequence of large $N$. An apparatus with $\sim\!10^{23}$ vertices has $\sim\!10^{23}$ phases relative to the measured system. Their sum:
\begin{equation}
\sum_{j=1}^{N_\text{app}} G_{sj} = \sum_{j} |G_{sj}|\,e^{i\phi_{sj}} \;\xrightarrow{N \to \infty}\; O(1/\sqrt{N}) \quad\text{(non-aligned phases cancel)}
\end{equation}
Only the component of $\psi_s$ aligned with the apparatus eigenstates survives. This is projection --- not as a postulate, but as the law of large numbers applied to phases.

\begin{remark}
The ``measurement problem'' does not arise. $G$ is static. Nothing collapses. The appearance of collapse is what reading a static matrix looks like when the reader is part of the matrix.
\end{remark}

\subsection{No observables, no wave-particle duality}

This theory derived 52 quantitative predictions --- coupling constants, fermion masses, mixing angles, cosmological observables, $\eta/s = 1/(4\pi)$, the proton mass to 0.08\% --- without defining a single observable operator. No position operator $\hat{x}$, no momentum $\hat{p}$, no Hamiltonian postulated as ``the energy operator,'' no eigenvalue-equals-measurement axiom, no Born rule. The local Hamiltonian $H_i = \sum_j W_{ij}|\psi_j\rangle\langle\psi_j|$ was not introduced --- it emerged from the Gram matrix structure.

The question ``is it a particle or a wave?'' does not arise. Each vertex carries a complex number $\psi \in \mathbb{C}^5$. Complex numbers are neither particles nor waves. The question is malformed --- like asking whether the number 3 is red or blue.

\subsection{The minimal assumption}

The results of this work follow from a single mathematical fact: \emph{complex numbers exist}. Given sufficiently many complex numbers $\{z_\alpha\}$ assigned to any combinatorial structure, one can compute the Gram matrix $G_{ij} = \sum_\alpha z^*_{i,\alpha}\,z_{j,\alpha}$. This matrix:
\begin{enumerate}
\item is Hermitian ($G^\dagger = G$),
\item is positive semidefinite ($\mathbf{c}^\dagger G\,\mathbf{c} \geq 0$),
\item has finite rank ($\text{rank}(G) \leq d$ where $d$ is the number of components per vertex),
\item stratifies by effective rank when sorted by $\det(G_h)$.
\end{enumerate}

For any $d$ large enough, the rank stratification contains a band at effective rank 5. In that band, the Binet--Cauchy decomposition gives $1 + 12 + 12 = 25$ channels, the Basel sum gives $1/\alpha_1 = 6\pi^2$, the causal split gives $\text{SU}(3) \times \text{SU}(2) \times \text{U}(1)$, and all 52 predictions follow.

No physics was assumed. No spacetime, no forces, no particles, no quantum mechanics, no relativity. Only $\mathbb{C}$ and the inner product.

\bigskip

\noindent\textbf{Theorem} (Existence of physics). \emph{Let $\{z_{i,a}\}$ be $N \times d$ complex numbers with $N \gg d \gg 5$, drawn from any distribution with finite variance. Let $G_{ij} = \sum_a z^*_{i,a}\,z_{j,a}$. Then with probability $1$, $G$ contains a region of effective rank $5$ whose geometric properties reproduce the Standard Model and general relativity.}

\bigskip

\noindent The proof is the content of this book.
